\bumppoints{22}

\question
The ``role'' or ``job'' of a species in a community, represented by all of the resources needed by the species, is called 

\begin{choices}
\choice a predator.
\choice a producer or consumer.
\correctchoice the niche.
\choice the resource.
\choice a predator or prey.
\end{choices}

\question
Which block in the figure below is most consistent with intermediate levels of disturbance in an ecosystem? (Different shades, including white, represent different species. The small squares indicate relative abundance of individuals present.)

\begin{multicols}{3}
	\begin{choices}
\choice \includegraphics{diversity_mcD} 
\choice \includegraphics{diversity_mcA} 
\choice \includegraphics{diversity_mcB} 
\choice \includegraphics{diversity_mcC} 
\correctchoice \includegraphics{diversity_mcE} 
\end{choices}
\end{multicols}

\question
In the experiment shown in the figure below, \textit{Balanus} was removed from the habitat shown on the left. Which of the following statements is a valid conclusion of this experiment?\\
\includegraphics[width=0.9\textwidth]{exam3_balanus}

\begin{choices}
\choice \textit{Balanus} can survive only in the lower intertidal zone because it is unable to resist dehydration.
\choice \textit{Balanus} is inferior to \textit{Chthamalus} in competing for space on rocks lower in the intertidal zone. 
\choice If \textit{Chthamalus} were removed, \textit{Balanus}'s fundamental niche would become larger.
\choice These two species of barnacle do not show competitive exclusion.
\correctchoice The removal of \textit{Balanus} shows that the realized niche of \textit{Chthamalus} is smaller than its fundamental niche.
\end{choices}

\newpage

\question
Cowbirds birds lay their eggs in the nests of other birds like cardinals. The cardinals raise the cowbird chicks as their own but the cowbird chicks usually kill the cardinal chicks of the other bird. This is most likely an example of 

\begin{choices}
\choice mutualism. 
\choice competition. 
\choice resource partitioning. 
\correctchoice parasitism. 
\choice symbiosis. 
\end{choices}

\question
Which of the following is \emph{not} directly concerned with community ecology?

\begin{choices}
\choice herbivory
\choice predation
\choice resource partitioning
\correctchoice intraspecific competition
\choice symbiosis
\end{choices}

\question
Several species of warblers (a group of small colorful birds) in North America can live in the same tree only because they

\begin{choices}
\choice can fly to different heights within the tree.
\choice out-competed other species of birds, such as sparrows.
\correctchoice partitioned resources to eliminate competition.
\choice eat different foods within the tree.
\choice created different types of nests within the tree.
\end{choices}

\question
Given the following values and assuming exponential population growth, what will be the total population size after one generation has elapsed? Round to nearest whole individual.\vspace{1ex}

\hspace{\leftmargin}%
\begin{minipage}{2in}
$b = 0.09$ \qquad $N = 18,657$\\
$d = 0.12$ \qquad $K = 15,000$\\
\end{minipage}

\begin{choices}
\choice $19,217$ 
\correctchoice $18,097$
\choice $560$
\choice $3,657$ 
\choice $15,000$
\end{choices}

\question
Competitive exclusion means that

\begin{choices}
\choice one individual will always out complete another individual in a population. 
\choice a species will have a smaller realized niche than a fundamental niche. 
\choice two species will partition resources to eliminate competition. 
\choice one population will always exclude another another population of the same species in a community. 
\correctchoice one species will always out complete another species that occupies the same niche.  
\end{choices}

\question
Color that helps hide an animal in its environment is called

\begin{choices}
\choice environmental coloration. 
\choice aposomatic coloration 
\correctchoice cryptic coloration. 
\choice hideandseek coloration. 
\choice mimicry. 
\choice warning coloration. 
\end{choices}

\question
The dispersion pattern for individuals in across the entire metapopulation shown in question \ref{ques:metapopulation} is most likely

\begin{choices}
\choice uniform. 
\choice even. 
\choice gene flow. 
\correctchoice clumped. 
\choice random. 
\end{choices}

\question
Which of the following is most likely to contribute to density-dependent regulation of populations?

\begin{choices}
\choice Earthquakes.
\choice The removal of toxic waste by decomposers.
\choice Fires.
\correctchoice Intraspecific competition for nutrients.
\choice Floods.
\end{choices}
